\section{Background and Related Work}

\subsection{Resource-constrained sensing}

Resource-constrained sensor scheduling has been studied extensively in estimation,
communication, and energy-aware monitoring. Estimation-oriented formulations
minimize state-estimation error or posterior covariance
\citep{Shi_2011,AlAhdab2025}. Energy-harvesting censoring policies optimize the
utility of transmitted measurements \citep{FernandezBes2015}, while
freshness-driven policies use age of information (AoI) or age of incorrect
information \citep{Kaul_2012,Jonah2026}. Reinforcement learning has also been
applied to tracking and belief-state scheduling \citep{Qu2022,Alali2024} and to
delay-aware remote estimation with broader scheduling objectives
\citep{Tran2026}. These objectives directly match evaluation criteria based on
measurement freshness or current uncertainty. The present forecasting problem
retains the resource constraints and uses forecast loss as the scheduling
objective. A channel receives greater scheduling value when its observations
reduce forecast loss over the evaluation horizon.

A forecast-loss objective is particularly relevant for heterogeneous sensing
systems. A meteorological core channel can provide inexpensive and broadly useful
measurements, while optional channels may provide their greatest forecasting
benefit near specific physical events. The validation-selected fixed-schedule
baseline therefore provides a meaningful comparator. The evaluation compares observation-dependent
scheduling with fixed and dynamic baselines and with a privileged true-label baseline.

\begin{table}[htbp]
  \centering
  \caption{Position of this work relative to nearby scheduling and sensing literatures.}
  \label{tab:related_work_positioning}
  \footnotesize
  \setlength{\tabcolsep}{3.0pt}
  \renewcommand{\arraystretch}{0.96}
  \begin{tabular}{@{}>{\raggedright\arraybackslash}p{0.23\linewidth}>{\raggedright\arraybackslash}p{0.25\linewidth}>{\raggedright\arraybackslash}p{0.22\linewidth}>{\raggedright\arraybackslash}p{0.20\linewidth}@{}}
    \toprule
    Literature line & Typical optimization target & Scheduling structure & What this paper adds \\
    \midrule
    State-estimation and AoI scheduling
      & Covariance, estimation error, information age
      & Sensor selection under communication or energy limits
      & Multi-step forecast loss is the objective. \\
    Active perception and adaptive observation selection
      & Task information under partial observability
      & Observation actions selected from partial observations
      & A mandatory meteorological core channel, limited optional-channel capacity,
        and baselines that use information unavailable during deployment. \\
    DRL-based adaptive sensing and sensor steering
      & Information gain, coverage, data acquisition utility, or monitoring reward
      & Learned adaptive policies, often with soft or task-specific constraints
      & A finite candidate action set enforces the mandatory-core and
        optional-channel-capacity rules, and a feasible-action mask checks the
        steady-state and startup-peak power limits. \\
    Time series forecasting models
      & Forecast accuracy after the observation stream is given
      & Forecasting model for a given observation stream
      & A frozen primary forecaster evaluates every schedule; its parameters are fixed during policy training and evaluation. \\
    \bottomrule
  \end{tabular}
\end{table}


\subsection{Adaptive sensing and learning-based scheduling}

Forecast-oriented sensor scheduling is related to active perception and partially
observable decision making. In these settings, an agent selects observations
sequentially to improve a downstream task
\citep{Bajcsy2018,Lauri_2023}. Related theoretical work studies adaptive
observation selection under partial observability \citep{Golovin2011}. These
formulations support policies that respond to available observations. They do not
address how forecast-oriented channel selections should be executed under
steady-state power limits, startup-peak power limits, and a minimum dwell-time
constraint.

Deep reinforcement learning has been used for resource-constrained adaptive
sensing and sensor steering \citep{Murad2020,Ogbodo2025}, informative path
planning \citep{Wei2020}, and constrained optimization with delayed rewards
\citep{Pendyala2024}. Its sequential objective can account for consequences that
emerge over multiple time steps. Explicit feasibility enforcement is required when
the sensing system cannot execute arbitrary actions. The PD-PPO policy is
categorical over candidate action indices, and the scheduler applies a
feasible-action mask at each time step. The mask removes indices whose candidate
activation vectors violate the steady-state or startup-peak power limit.

\subsection{Forecasting models for schedule evaluation}

Forecasting models can quantify the downstream prediction loss associated with
an observation sequence produced by a sensor schedule. Relevant architectures
include the Temporal Fusion Transformer, temporal convolutional networks, and
iTransformer \citep{Lim2021,Bai2018,Liu2024}. The primary forecaster is trained
on the forecaster-training partition, and its parameters are held fixed throughout
policy training and evaluation. During policy training, it provides the
multi-step forecast-loss term used to construct the policy reward. During test
evaluation, it measures the forecast loss of the observation sequence produced
by each schedule. The alternative-forecaster analysis applies a ridge forecaster
to the same observation sequences and reselects the fixed-schedule baseline on the
calibration/validation partition. This sensitivity analysis quantifies the same
comparison under a second forecasting model.

\subsection{Antarctic blowing-snow monitoring}

Antarctic automatic weather station (AWS) networks and blowing-snow monitoring
provide the application context for the simulation. Antarctic AWS networks supply long-term
atmospheric and surface observations, while their spatial instrument coverage
remains uneven \citep{AntAWS2023}. Field measurements and numerical models
have been used to study drifting-snow occurrence and transport,
atmosphere--snow interactions, and snow particle properties
\citep{Amory2020,Sharma2023,Monrad2026}. The controlled simulation represents
this setting with one mandatory meteorological core channel and five optional channels.
At most one optional channel can operate at a time. For each seed, the PD-PPO scheduler and all
baselines are evaluated on the same simulated ground-truth time series.
